\documentclass[../../../include/open-logic-section]{subfiles}

\begin{document}

\olfileid{sth}{replacement}{strength}
\olsection{قوة الاستبدال}

نستفتح بمثال يسير يكشف قوة الاستبدال: ما دمنا لا نتجاوز~$\Z$، فلا
سبيل إلى إثبات وجود عدد ترتيبي لفون نويمان يبلغ~$\omega + \omega$
أو يجاوزه.

وبيان وجه ذلك على سبيل الرسم كما يلي. نعمل في~$\ZF$ وننظر إلى
$\OLId{latin-looped}{V}_{\omega+\omega}$؛ فإن هذه المجموعة تصلح مجالًا لـ\emph{نموذج}
من~$\Z$. ولإيضاح هذا ندخل ترميز نسبنة الصيغ.
\begin{defn}\ollabel{formularelativization}
	لكل مجموعة~$\OLId{latin-looped}{M}$ وصيغة~$\phi$، نرمز بـ$\phi^\OLId{latin-looped}{M}$ إلى الصيغة الناتجة
	من تقييد كل مكمم في~$\phi$ بـ$\OLId{latin-looped}{M}$؛ أي نستبدل «$\lexists[\OLId{latin-ordinary}{x}]$» بـ
	«$(\lexists[\OLId{latin-ordinary}{x} \in \OLId{latin-looped}{M}])$»، ونستبدل «$\lforall[\OLId{latin-ordinary}{x}]$» بـ
	«$(\lforall[\OLId{latin-ordinary}{x} \in \OLId{latin-looped}{M}])$».
\end{defn}\noindent
ويمكن البرهنة على أن كل مسلّمة~$\phi$ من مسلّمات~$\Z$ تحقق
$\ZF \vdash \phi^{\OLId{latin-looped}{V}_{\omega+\omega}}$. غير أن~$\omega+\omega$ ليست
\emph{في}~$\OLId{latin-looped}{V}_{\omega+\omega}$، بمقتضى
\olref[spine][rank]{ordsetrankalpha}. فيتسق~$\Z$، من ثم، مع عدم وجود
$\omega+\omega$.

وهنا يظهر سبب قولنا في~\olref[ordinals][replacement]{sec} إن
\olref[ordinals][ordtype]{thmOrdinalRepresentation} لا يثبت بغير
الاستبدال. فنستطيع داخل~$\Z$ أن نعرّف صراحة ترتيبًا جيدًا \emph{حقه}
في الحدس أن يكون نمط ترتيبه~$\omega+\omega$. وقد مر مثال غير صوري في
\olref[ordinals][idea]{sec}، هو ترتيب الأعداد الطبيعية المحدد بـ:
\begin{align*}
	\OLId{latin-ordinary}{n} \lessdot \OLId{latin-ordinary}{m} \text{ إذا وفقط إذا كان }&
	\text{إما }\OLId{latin-ordinary}{n} < \OLId{latin-ordinary}{m}\text{ و }\OLId{latin-ordinary}{m}-\OLId{latin-ordinary}{n}\text{ زوجيًا،}\\
	& \text{أو كان $n$ زوجيًا و$m$ فرديًا.}
\end{align*}
فإذا لم يوجد~$\omega+\omega$، لم يتماثل هذا الترتيب الجيد مع عدد
ترتيبي قط. ولذلك لا تثبت~$\Z$
\olref[ordinals][ordtype]{thmOrdinalRepresentation}.

وإذا قلبنا النظر رأينا أن الاستبدال يثبت وجود~$\omega+\omega$، فيلزم
أن يثبت وجود~$\OLId{latin-looped}{V}_{\omega+\omega}$ أيضًا. بل الأمر أعم: فكل ترتيب جيد
نقدر على تعريفه يقابله، بحكم
\olref[ordinals][ordtype]{thmOrdinalRepresentation}، عدد ترتيبي
$\OLId{greek-variable}{alpha}$ يتماثل معه؛ ومن ثم يوجد~$\OLId{latin-looped}{V}_\OLId{greek-variable}{alpha}$. فالاستبدال هو الذي يكفل،
على وجه مباشر، أن يمتد هرم المجموعات إلى \emph{ارتفاع عظيم}.

وفي الأقسام التالية، ثم في~\olref[card-arithmetic][fix]{sec}، يتضح
مقدار ذلك الارتفاع أكثر. وحسبنا هنا أن قوة الاستبدال نفسها تجعل طلب
تسويغه أمرًا لا محيد عنه.

\end{document}
